\subsection{Interpretation of the Hubble Tension as an Environmental Effect}
  \label{subsec:interpretation-of-the-hubble-tension-as-an-environmental-effect}

  The framework developed here suggests that the Hubble tension does not reflect a
  change in the global expansion history.
  It reflects a limitation of homogeneous kinematic inference when applied to a
  strongly inhomogeneous late-time Universe.

  Early-universe determinations of \(H_0\),
  anchored in a nearly homogeneous background,
  recover the global expansion rate.
  By contrast, late-time local measurements probe structured environments in which
  the effective transition resolution operates in a low-density saturated regime.

  In this context, the discrepancy between early- and late-time determinations of
  \(H_0\) arises from a mismatch between global averaging and local inference.
  Measurements performed in void-dominated regions sample an effective response
  that no longer scales linearly with decreasing density,
  leading to a systematically enhanced locally inferred expansion rate.

  From this perspective, the Hubble tension is not an anomaly requiring additional
  dark components or early-time modifications.
  It is a diagnostic of gravitational and kinematic inference in diffuse
  environments.
  This interpretation naturally predicts weak but systematic correlations between
  inferred values of \(H_0\) and large-scale environmental indicators such as void
  depth or filament membership.
